\documentclass[12pt,a4paper]{article}

% ── Paquetes ──────────────────────────────────────────────────────────────────
\usepackage[utf8]{inputenc}
\usepackage[T1]{fontenc}

\usepackage{amsmath,amssymb,amsthm}
\usepackage{graphicx}
\usepackage{booktabs}
\usepackage{array}
\usepackage{caption}
\usepackage{subcaption}
\usepackage{hyperref}
\usepackage{geometry}
\usepackage{fancyhdr}
\usepackage{setspace}
\usepackage{xcolor}
\usepackage{listings}
\usepackage{float}



% ── Geometría ─────────────────────────────────────────────────────────────────
\setlength{\headheight}{14pt}
\geometry{
  a4paper,
  top=2.5cm, bottom=2.5cm,
  left=2.5cm, right=2.5cm
}

% ── Encabezado y pie ──────────────────────────────────────────────────────────
\pagestyle{fancy}
\fancyhf{}
\rhead{\small Análisis numérico de EDPs · Proyecto Final}
\lhead{\small Pista B: Problema Inverso con PINNs}
\cfoot{\thepage}
\renewcommand{\headrulewidth}{0.4pt}

% ── Hipervínculos ─────────────────────────────────────────────────────────────
\hypersetup{
  colorlinks=true,
  linkcolor=blue!70!black,
  citecolor=blue!70!black,
  urlcolor=blue!70!black
}

% ── Interlineado ──────────────────────────────────────────────────────────────
\onehalfspacing

% ── Macros ────────────────────────────────────────────────────────────────────
\newcommand{\lhat}{\hat{\lambda}}
\newcommand{\ltrue}{\lambda_{\text{real}}}
\newcommand{\Lf}{\mathcal{L}_f}
\newcommand{\Lb}{\mathcal{L}_b}
\newcommand{\Li}{\mathcal{L}_{ic}}
\newcommand{\Ld}{\mathcal{L}_d}

% ══════════════════════════════════════════════════════════════════════════════
\begin{document}

% ── Portada ───────────────────────────────────────────────────────────────────
\begin{titlepage}
  \centering
  \vspace*{2cm}
  {\Large \textbf{Análisis Numérico de Ecuaciones Diferenciales} \par}
  \vspace{0.5cm}
  {\large Proyecto Final \par}
  \vspace{2cm}
  {\LARGE \textbf{Redes Neuronales Informadas por la Física (PINNs):} \par}
  \vspace{0.4cm}
  {\LARGE \textbf{Problema Inverso de Parámetros} \par}
  \vspace{0.5cm}
  {\large Pista B \par}
  \vspace{3cm}
  {\large \textit{Ecuación de difusión 1D con coeficiente desconocido} \par}
  \vspace{3cm}
  {\normalsize Basado en: Raissi, Perdikaris y Karniadakis (2019),\\
  \textit{J. Comput. Phys.} 378:686--707 \par}
  \vfill
  {\normalsize \today \par}
\end{titlepage}

% ── Tabla de contenidos ───────────────────────────────────────────────────────
\tableofcontents
\newpage

% ══════════════════════════════════════════════════════════════════════════════
\section{Introducción y Planteamiento del Problema}

En numerosos fenómenos de transporte, la forma funcional de la ecuación diferencial en
derivadas parciales (EDP) es conocida, pero ciertos coeficientes constitutivos ---como la
difusividad térmica o la viscosidad--- permanecen inaccesibles a medición directa. El
\textbf{problema inverso de parámetros} consiste en inferir dichos coeficientes a partir de
observaciones dispersas y ruidosas del campo solución, sin recurrir a experimentos costosos
ni a discretizaciones globales del dominio.

Se estudia la \textbf{ecuación de difusión unidimensional}:
\begin{equation}
  \frac{\partial u}{\partial t} = \lambda\,\frac{\partial^2 u}{\partial x^2},
  \qquad (x,t)\in[0,1]^2,
  \label{eq:difusion}
\end{equation}
con condición inicial $u(x,0)=\sin(\pi x)$ y condiciones de contorno Dirichlet homogéneas
$u(0,t)=u(1,t)=0$. La solución analítica exacta es
\begin{equation}
  u(x,t)=\sin(\pi x)\,e^{-\lambda\pi^2 t},
  \label{eq:analitica}
\end{equation}
lo que permite generar datos sintéticos de referencia y evaluar rigurosamente los errores.
El parámetro de difusividad $\lambda>0$ se trata como \textbf{incógnita}; el objetivo es
recuperarlo conjuntamente con el campo completo $u(x,t)$ a partir de $N_d=100$
observaciones distribuidas aleatoriamente en el dominio, contaminadas con un 1\,\% de
ruido gaussiano relativo.

% ══════════════════════════════════════════════════════════════════════════════
\section{Metodología}

\subsection{Arquitectura de la red}

Se implementó una PINN (\textit{Physics-Informed Neural Network}) como subclase de
\texttt{tf.keras.Model} en TensorFlow 2.x, con la siguiente arquitectura:
\begin{itemize}
  \item \textbf{Entradas:} coordenadas $(x,t)\in\mathbb{R}^2$, concatenadas antes de la
        primera capa oculta.
  \item \textbf{Capas ocultas:} 3 capas densas de 40 neuronas cada una, con activación
        $\tanh$.
  \item \textbf{Salida:} aproximación escalar $\tilde{u}(x,t;\theta)$.
  \item \textbf{Inicialización:} pesos con distribución Glorot normal, adecuada para
        activaciones $\tanh$.
\end{itemize}

El coeficiente desconocido se declara como \texttt{tf.Variable} escalar independiente de
los pesos de la red, inicializado en $\lhat_0=0.5$ (cinco veces el valor real), con
restricción explícita de positividad $\lhat\geq10^{-6}$ aplicada tras cada paso de
gradiente.

\subsection{Función de pérdida}

La pérdida total combina cuatro términos:
\begin{equation}
  \mathcal{L}(\theta,\lhat)
  = w_f\,\Lf + \Li + \Lb + w_d\,\Ld,
  \label{eq:loss}
\end{equation}
donde $w_f=1$ y $w_d=10$, y cada término es:
\begin{align}
  \Lf &= \frac{1}{N_f}\sum_{i=1}^{N_f}
         \left(\frac{\partial\tilde{u}}{\partial t}\bigg|_i
         -\lhat\,\frac{\partial^2\tilde{u}}{\partial x^2}\bigg|_i\right)^2, \\[4pt]
  \Li &= \frac{1}{N_{ic}}\sum_{i=1}^{N_{ic}}
         \bigl(\tilde{u}(x_i,0)-\sin(\pi x_i)\bigr)^2, \\[4pt]
  \Lb &= \frac{1}{N_{bc}}\sum_{i=1}^{N_{bc}}
         \bigl(\tilde{u}(0,t_i)^2+\tilde{u}(1,t_i)^2\bigr), \\[4pt]
  \Ld &= \frac{1}{N_d}\sum_{i=1}^{N_d}
         \bigl(\tilde{u}(x_i^d,t_i^d)-u_i^{\text{obs}}\bigr)^2.
\end{align}

Se emplearon $N_f=2000$ puntos de colocación y $N_{ic}=N_{bc}=100$.

\subsection{Diferenciación automática}

Las derivadas parciales $\partial\tilde{u}/\partial t$ y
$\partial^2\tilde{u}/\partial x^2$ se obtienen mediante diferenciación automática con
\texttt{tf.GradientTape}. Se utilizaron \textbf{tapes anidados} para garantizar el cálculo
correcto de la segunda derivada:
\begin{itemize}
  \item El tape externo (\texttt{tape2}) traza $\tilde{u}$ y $\partial\tilde{u}/\partial x$
        respecto a $x$.
  \item El tape interno (\texttt{tape1}) calcula $\partial\tilde{u}/\partial t$ y
        $\partial\tilde{u}/\partial x$.
  \item La segunda derivada $\partial^2\tilde{u}/\partial x^2$ se obtiene aplicando
        \texttt{tape2} sobre $\partial\tilde{u}/\partial x$.
\end{itemize}
Este diseño evita el error sutil de reutilizar un tape persistente para la segunda
derivada, donde la primera derivada no habría sido trazada en el grafo del tape exterior.

\subsection{Entrenamiento}

Las variables entrenables se componen de los parámetros de la red $\theta$ más el escalar
$\lhat$, optimizados \textbf{conjuntamente} mediante Adam con tasa de aprendizaje
$\eta=10^{-3}$ durante 8000 épocas. El gradiente $|\nabla_{\lhat}\mathcal{L}|$ se registró
en cada iteración como indicador de diagnóstico, permitiendo detectar estancamiento o
colapso del parámetro en etapas tempranas del entrenamiento.

% ── Figura: condición inicial ──────────────────────────────────────────────────
\begin{figure}[H]
  \centering
  \includegraphics[width=0.5\textwidth]{13.png}
  \caption{Condición inicial $u(x,0)=\sin(\pi x)$ utilizada en todos los experimentos.}
  \label{fig:ic}
\end{figure}

% ── Figura: datos observados ───────────────────────────────────────────────────
\begin{figure}[H]
  \centering
  \includegraphics[width=0.5\textwidth]{10.png}
  \caption{Distribución de los $N_d=100$ puntos de observación en el dominio
           $(x,t)\in[0,1]^2$, coloreados por el valor de $u$ observado con 1\,\%
           de ruido gaussiano.}
  \label{fig:datos}
\end{figure}

% ══════════════════════════════════════════════════════════════════════════════
\section{Problema Directo}

Como caso de referencia, se resolvió primero el \textbf{problema directo} con
$\lambda=0.1$ conocido, entrenando el PINN únicamente con los términos $\Lf$, $\Li$ y
$\Lb$. Este experimento valida la capacidad de la red para aproximar la solución de la EDP
antes de añadir la complejidad del problema inverso.

\subsection{Convergencia}

La Figura~\ref{fig:conv_directo} muestra la evolución de la pérdida total y sus
componentes durante el entrenamiento. Las tres pérdidas decrecen de forma monótona, con la
pérdida total alcanzando $\sim10^{-4}$ al finalizar las épocas de entrenamiento.

\begin{figure}[H]
  \centering
  \includegraphics[width=0.8\textwidth]{12.png}
  \caption{Convergencia del problema directo: pérdida total, residuo PDE y
           condiciones de contorno/iniciales en escala logarítmica.}
  \label{fig:conv_directo}
\end{figure}

\subsection{Solución reconstruida}

La Figura~\ref{fig:sol_directo} compara la solución analítica con la aproximación del
PINN en el dominio completo, junto con el mapa de error absoluto. El error se concentra
en la región $t\to1$, donde la solución es pequeña y el residuo relativo crece.

\begin{figure}[H]
  \centering
  \includegraphics[width=\textwidth]{11.png}
  \caption{Problema directo: solución analítica (izquierda), solución PINN (centro) y
           error absoluto (derecha) en el dominio $[0,1]^2$.}
  \label{fig:sol_directo}
\end{figure}

% ══════════════════════════════════════════════════════════════════════════════
\section{Problema Inverso}

\subsection{Convergencia del parámetro}

La Figura~\ref{fig:conv_lambda} muestra la evolución conjunta de $\lhat$ y la pérdida
total durante el entrenamiento. El parámetro converge desde $\lhat_0=0.5$ hasta el
vecindario del valor real $\ltrue=0.1$ en menos de 1000 épocas, y se estabiliza a
partir de la época 2500.

\begin{figure}[H]
  \centering
  \includegraphics[width=0.75\textwidth]{9.png}
  \caption{Evolución de $\lhat$ durante el entrenamiento del problema inverso.
           La línea roja discontinua indica el valor real $\ltrue=0.1$.}
  \label{fig:conv_lambda}
\end{figure}

\begin{figure}[H]
  \centering
  \includegraphics[width=0.75\textwidth]{8.png}
  \caption{Evolución de la pérdida total durante el entrenamiento del problema inverso
           en escala logarítmica.}
  \label{fig:loss_inv}
\end{figure}

\subsection{Resultado final}

Al finalizar las 8000 épocas de entrenamiento, los resultados obtenidos fueron:

\begin{table}[H]
  \centering
  \caption{Resultado del problema inverso.}
  \label{tab:resultado}
  \begin{tabular}{lc}
    \toprule
    Magnitud & Valor \\
    \midrule
    $\ltrue$              & 0.10000 \\
    $\lhat$ estimado      & 0.10017 \\
    Error relativo        & \textbf{0.17\,\%} \\
    $|\nabla_{\lhat}\mathcal{L}|$ final & $1.58\times10^{-6}$ \\
    \bottomrule
  \end{tabular}
\end{table}

El error relativo de 0.17\,\% supera con amplio margen el umbral de referencia del 1\,\%
establecido en el enunciado.

\subsection{Reconstrucción del campo}

La Figura~\ref{fig:sol_inverso} muestra la comparación entre la solución PINN y la
analítica evaluadas con $\lhat=0.10017$, junto con el mapa de error absoluto. El patrón
de error es similar al del problema directo, con máximos en la región $t\to1$,
$x\to\{0,1\}$.

\begin{figure}[H]
  \centering
  \includegraphics[width=\textwidth]{4.png}
  \caption{Problema inverso: solución PINN con $\lhat=0.10017$ (izquierda), solución
           analítica (centro) y error absoluto (derecha).}
  \label{fig:sol_inverso}
\end{figure}

% ══════════════════════════════════════════════════════════════════════════════
\section{Validación}

\subsection{Comparación con diferencias finitas (Crank-Nicolson)}

Como primer control de validación, se implementó un solver de referencia basado en el
esquema de \textbf{Crank-Nicolson} (CN), incondicionalmente estable, con $N_x=99$ nodos
interiores y $N_t=1000$ pasos temporales. La comparación se realizó evaluando los tres
métodos ---solución analítica, CN y PINN--- sobre la misma malla, reportando el error
$L^2$ relativo $\|u_A-u_B\|_2/\|u_A\|_2$.

\begin{table}[H]
  \centering
  \caption{Errores $L^2$ relativos de la validación contra Crank-Nicolson.}
  \label{tab:cn}
  \begin{tabular}{lcc}
    \toprule
    Comparación & Problema Directo & Problema Inverso ($\lhat=0.10017$) \\
    \midrule
    CN vs.\ analítica   & $6.711\times10^{-3}$ & $6.737\times10^{-3}$ \\
    PINN vs.\ analítica & $1.470\times10^{-3}$ & $1.565\times10^{-3}$ \\
    CN vs.\ PINN        & $7.453\times10^{-3}$ & $7.250\times10^{-3}$ \\
    \bottomrule
  \end{tabular}
\end{table}

El PINN supera en precisión al esquema CN respecto a la solución analítica en ambos casos.
Esta diferencia se explica por el origen del error en cada método: el CN acumula error de
truncación de orden $\mathcal{O}(\Delta x^2,\Delta t^2)$, mientras que el PINN minimiza
directamente el residuo continuo de la EDP. La concordancia mutua ---diferencias del
orden de $10^{-3}$--- confirma que ambos métodos son consistentes y que la solución del
PINN es físicamente correcta.

\begin{figure}[H]
  \centering
  \includegraphics[width=\textwidth]{3.png}
  \caption{Problema directo: cortes temporales en $t\in\{0,0.25,0.5,0.75,1.0\}$
           para la solución analítica, Crank-Nicolson y PINN.}
  \label{fig:cn_directo}
\end{figure}

\begin{figure}[H]
  \centering
  \includegraphics[width=\textwidth]{2.png}
  \caption{Problema inverso: cortes temporales en $t\in\{0,0.25,0.5,0.75,1.0\}$
           para la solución analítica, Crank-Nicolson y PINN ($\lhat=0.10017$).}
  \label{fig:cn_inverso}
\end{figure}

\subsection{Estudio de sensibilidad}

Para evaluar la robustez del método, se varió el número de puntos de observación
$N_d\in\{20,50,100,200\}$ y el nivel de ruido $\sigma\in\{0\%,1\%\}$, entrenando cada
configuración durante 4000 épocas. El experimento se ejecutó con \textbf{tres semillas
independientes} para evaluar la variabilidad estocástica.

\begin{table}[H]
  \centering
  \caption{Error relativo en $\lhat$ (\%) para distintos $N_d$ y niveles de ruido
           (media sobre tres semillas, con rango entre paréntesis).}
  \label{tab:sensibilidad}
  \begin{tabular}{ccc}
    \toprule
    $N_d$ & Ruido 0\,\% & Ruido 1\,\% \\
    \midrule
    20  & 0.18\,\% (0.03--0.26) & 0.35\,\% (0.11--0.52) \\
    50  & 0.23\,\% (0.17--0.29) & 0.12\,\% (0.07--0.21) \\
    100 & 0.21\,\% (0.10--0.29) & 0.65\,\% (0.53--0.80) \\
    200 & 0.19\,\% (0.14--0.28) & 0.26\,\% (0.07--0.65) \\
    \bottomrule
  \end{tabular}
\end{table}

En la totalidad de las 24 configuraciones evaluadas, el error en $\lhat$ permaneció
\textbf{por debajo del 1\,\%}, confirmando que el método recupera el coeficiente de
difusión de forma robusta incluso con tan solo 20 observaciones. La Figura~\ref{fig:sens}
muestra la evolución del error en $\lhat$ y del error $L^2$ del campo en función de $N_d$.

\begin{figure}[H]
  \centering
  \includegraphics[width=\textwidth]{1.png}
  \caption{Estudio de sensibilidad: error relativo en $\lhat$ (izquierda) y error $L^2$
           del campo reconstruido (derecha) en función de $N_d$, para dos niveles de
           ruido. La línea discontinua gris indica el umbral de referencia del 1\,\%.}
  \label{fig:sens}
\end{figure}

\subsection{Validación de plausibilidad física}

Se realizaron los siguientes controles de plausibilidad:
\begin{enumerate}
  \item \textbf{Positividad del parámetro:} se garantizó $\lhat>0$ en todo el
        entrenamiento mediante proyección explícita $\lhat\leftarrow
        \max(\lhat,10^{-6})$ tras cada actualización de gradiente.
  \item \textbf{Consistencia del campo reconstruido:} la solución $\tilde{u}(x,t;\theta^*)$
        evaluada con $\lhat^*$ reproduce fielmente los contornos y la dinámica de
        decaimiento exponencial de la solución analítica.
  \item \textbf{Diagnóstico de gradiente:} el seguimiento de $|\nabla_{\lhat}\mathcal{L}|$
        a lo largo del entrenamiento confirmó que el parámetro recibió señal de gradiente
        consistente y no colapsó hacia el límite inferior.
\end{enumerate}

% ══════════════════════════════════════════════════════════════════════════════
\section{Discusión y Limitaciones}

\subsubsection*{Paisaje de optimización no convexo.}
La función de pérdida $\mathcal{L}(\theta,\lhat)$ es altamente no convexa. El estudio de
sensibilidad reveló un caso anómalo aislado ($N_d=50$, ruido 0\,\%, primera semilla:
error $L^2=4.12\times10^{-2}$) no reproducido en las dos corridas posteriores, que se
atribuye a un mínimo local espurio del optimizador. Esto motiva el uso de múltiples
inicializaciones en aplicaciones críticas.

\subsubsection*{Ausencia de tendencia monótona con $N_d$.}
Para ruido nulo, el error en $\lhat$ no mejora monótonamente al aumentar $N_d$, oscilando
entre 0.18\,\% y 0.23\,\%. Esto indica que el factor limitante es el paisaje de
optimización, no la cantidad de datos, al menos en el rango explorado.

\subsubsection*{Variabilidad con ruido.}
Con ruido del 1\,\%, la variabilidad inter-semilla es mayor, especialmente en $N_d=100$
(rango de 0.27 puntos porcentuales). Este comportamiento es consistente con la mayor
dificultad del problema cuando los datos son ruidosos y la distribución aleatoria de los
puntos interactúa con el paisaje de pérdida.

\subsubsection*{Extensión a Fokker-Planck.}
Una extensión natural del método es su aplicación a la ecuación de Fokker-Planck para
inferir conjuntamente los términos de deriva y difusión a partir de observaciones de la
densidad de probabilidad, siguiendo la línea propuesta por Chen et al.\ (2021).

% ══════════════════════════════════════════════════════════════════════════════
\section{Conclusiones}

Se implementó y validó exitosamente una PINN para resolver el problema inverso de
parámetros en la ecuación de difusión 1D. Los resultados principales son:

\begin{enumerate}
  \item El coeficiente de difusividad $\lambda$ fue recuperado con un error relativo del
        \textbf{0.17\,\%}, superando con margen el umbral del 1\,\% establecido en el
        enunciado.
  \item La reconstrucción del campo $u(x,t)$ alcanzó un error $L^2$ relativo de
        $1.57\times10^{-3}$, consistente con el objetivo de $\sim10^{-3}$.
  \item La validación contra Crank-Nicolson confirmó la consistencia física de la
        solución; el PINN superó en precisión al esquema clásico con $N_x=99$ y
        $N_t=1000$.
  \item El estudio de sensibilidad demostró que el método es robusto para
        $N_d\geq20$ y niveles de ruido de hasta el 1\,\%, con todas las configuraciones
        evaluadas por debajo del umbral de referencia.
\end{enumerate}

% ── Referencias ───────────────────────────────────────────────────────────────
\bibliographystyle{plain}
\begin{thebibliography}{9}

\bibitem{raissi2019}
Raissi, M., Perdikaris, P., \& Karniadakis, G.\,E. (2019).
Physics-informed neural networks: a deep learning framework for solving forward and
inverse problems involving nonlinear partial differential equations.
\textit{Journal of Computational Physics}, 378, 686--707.

\bibitem{karniadakis2021}
Karniadakis, G.\,E., Kevrekidis, I.\,G., Lu, L., Perdikaris, P., Wang, S., \& Yang, L.
(2021). Physics-informed machine learning.
\textit{Nature Reviews Physics}, 3, 422--440.

\bibitem{chen2021}
Chen, X. et al.\ (2021). Solving inverse stochastic problems from discrete particle
observations using the Fokker-Planck equation and physics-informed neural networks.
\textit{SIAM Journal on Scientific Computing}.

\bibitem{krishnapriyan2021}
Krishnapriyan, A.\,S. et al.\ (2021). Characterizing possible failure modes in
physics-informed neural networks.
\textit{Advances in Neural Information Processing Systems}, 34.

\end{thebibliography}

\end{document}
